% Figura 1.2: Problemas de agencia y mecanismos de gobierno corporativo (Tema 1 UGR y Jensen-Meckling 1976)
\begin{tikzpicture}[
    node distance=1.3cm,
    principal/.style={
        rectangle, rounded corners=4pt, draw=dfinnavy, fill=dfinnavy!8,
        line width=1.1pt, minimum width=3.15cm, minimum height=1.55cm,
        align=center, font=\footnotesize\sffamily, text=dfinnavy
    },
    agente/.style={
        rectangle, rounded corners=4pt, draw=dfincopper, fill=dfincopper!10,
        line width=1.1pt, minimum width=3.2cm, minimum height=1.55cm,
        align=center, font=\footnotesize\sffamily, text=dfincopper
    },
    acreedor/.style={
        rectangle, rounded corners=4pt, draw=dfinemerald, fill=dfinemerald!8,
        line width=1.1pt, minimum width=3.15cm, minimum height=1.55cm,
        align=center, font=\footnotesize\sffamily, text=dfinemerald
    },
    mecanismo/.style={
        rectangle, rounded corners=3pt, draw=dfinslate!60, fill=dfinpanelivory,
        line width=0.8pt, text width=12.4cm, align=center,
        font=\scriptsize\sffamily, inner sep=6pt, text=dfinslate
    },
    rel/.style={<->, >=Latex, line width=1pt, color=#1},
    lbl/.style={font=\scriptsize\sffamily, align=center, text=#1}
]

\node[principal] (acc) at (-0.6, 2.55) {%
    \textbf{Accionistas (Principal)}\\[2pt]
    {\scriptsize\color{dfinslate}Objetivo: maximizar riqueza}\\
    {\scriptsize\color{dfinslate}($P_0$ y valor de mercado $E$)}
};

\node[agente] (dir) at (6.0, 2.55) {%
    \textbf{Directivos (Agente)}\\[2pt]
    {\scriptsize\color{dfinslate}Objetivo: salario, poder,}\\
    {\scriptsize\color{dfinslate}bajo esfuerzo y consolidación}
};

\node[acreedor] (acr) at (12.6, 2.55) {%
    \textbf{Acreedores (Deuda)}\\[2pt]
    {\scriptsize\color{dfinslate}Objetivo: cobro íntegro de}\\
    {\scriptsize\color{dfinslate}intereses y principal ($D$)}
};

% Conflicto 1: Accionistas <-> Directivos
\draw[rel=dfinburg] ([yshift=0.32cm]acc.east) -- ([yshift=0.32cm]dir.west)
    node[midway, lbl=dfinburg, above=3pt] {Construcción de imperios\\[-1pt]y aversión al riesgo};

% Conflicto 2: Directivos/Accionistas <-> Acreedores
\draw[rel=dfincopper] ([yshift=0.32cm]dir.east) -- ([yshift=0.32cm]acr.west)
    node[midway, lbl=dfincopper, above=3pt] {Sustitución de activos\\[-1pt]y dilución de deuda};

% Costes de agencia abajo
\draw[dashed, color=dfinmuted, line width=0.8pt] ([yshift=-0.32cm]acc.east) -- ([yshift=-0.32cm]dir.west)
    node[midway, lbl=dfinmuted, below=3pt] {Supervisión + fianza\\[-1pt]+ pérdida residual};

\draw[dashed, color=dfinmuted, line width=0.8pt] ([yshift=-0.32cm]dir.east) -- ([yshift=-0.32cm]acr.west)
    node[midway, lbl=dfinmuted, below=3pt] {\emph{Covenants}\\[-1pt]+ prima de crédito};

% Franja inferior de mecanismos alineadores
\node[mecanismo] (gov) at (6.0, -0.35) {%
    \textbf{\color{dfinnavy}Mecanismos de alineación y gobierno corporativo}\\[3pt]
    \textbf{Internos:} Consejo de administración independiente, auditoría, remuneración vinculada a acciones/TSR, deuda disciplinante.\\
    \textbf{Externos:} Mercado de control corporativo (OPAs hostiles), analistas financieros, agencias de rating y activismo institucional.
};

\draw[->, >=Latex, line width=0.9pt, color=dfinnavy] (gov.north -| acc) -- (acc.south);
\draw[->, >=Latex, line width=0.9pt, color=dfinnavy] (gov.north) -- (dir.south);
\draw[->, >=Latex, line width=0.9pt, color=dfinnavy] (gov.north -| acr) -- (acr.south);

\end{tikzpicture}
